\documentclass[10pt,a4paper]{article}
\usepackage[margin=19mm]{geometry}
\usepackage{amsmath,amssymb,amsthm,booktabs,microtype}
\usepackage[colorlinks=true,urlcolor=blue,citecolor=blue,linkcolor=blue]{hyperref}
\setlength{\parindent}{0pt}
\setlength{\parskip}{4pt}
\setlength{\emergencystretch}{2em}
\newcommand{\tht}{\vartheta}
\newcommand{\kap}{\kappa}
\newcommand{\gain}{G_\nu}
\newtheorem{theorem}{Theorem}
\newtheorem{lemma}{Lemma}
\title{Balanced source energies and adaptive marked moments}
\author{Zhongpai Gao}
\date{17 October 2026}
\begin{document}
\maketitle
\vspace{-2em}
\begin{abstract}
We retain a geometric surplus inside the source energies used in the
fixed-strip method for Dirichlet $L$-functions. Combined with
density-dependent inverse and plain cutoffs, this preserves the full
asymmetric marked gain through the allocation of prime factors.
We prove two uniform affine source estimates and combine them in a
finite decreasing count envelope. Using different moments for counting
and for the physical numerator gives the nonvanishing boundary
$0.874830$ for every Dirichlet $L$-function, while retaining the
established inverse-count exponent $5/6$.
\end{abstract}

\section{The reduction and the geometric saving}
The fixed-strip framework of OpenAI and Collet \cite{openai,collet}
uses source moments twice: to bound the physical numerator and to count
the detector's exceptional rows. The two uses need not select the same
moment. We exploit this flexibility after refining the source normalization.

Normalize logarithmic lengths by $\log U$, $U>1$. Let $m_1,m_2$ be the
two plain lengths and $s=\sum_j\log\mathrm N(P_j)/\log U$ the selected
prime width. Let $\mathcal E_\eta(\Phi,\mathbf f)$ be the inherited
initial affine energy, with its original masks, natural zeros, selected
weights, two profiles and independent heights bounded by $T$.
\begin{theorem}[Uniform affine sources]\label{thm:moment}
For each pair
\[
 (B_A,\lambda_A)=(0.9816147,3.7961251),\qquad
 (B_C,\lambda_C)=(0.97899425,3.72765011),
\]
the original source satisfies, for every $\varepsilon>0$,
\begin{equation}\label{eq:moment}
 \mathcal E_\eta(\Phi,\mathbf f)\ll_{\eta,\varepsilon}
 \mathcal D(\Phi)\mathcal S(\mathbf f)^2(1+T)^N U^{1+\varepsilon}
 \quad\text{if }m_1+m_2+\lambda_Ps\le B_P,
\end{equation}
with $0.405\le m_i\le0.44$ and $0.034\le s\le0.04$.
For fixed ray data and loss, a common positive slot mesh precedes the
ambient dictionary and slot cardinality. For each then fixed finite
family, the degree $N$, controls and constants precede moving sources,
subsets, row characters and heights, as in the initial-moment theorem.
\end{theorem}
Here $\mathcal D$ is the original diagonal control and $\mathcal S$ a
finite Schwartz control. Every displayed terminating decimal is exact.
\begin{theorem}[Dirichlet nonvanishing]\label{thm:main}
Every Dirichlet character satisfies
\[
 L(\chi,z)\ne0\quad\text{when }\Re z>\tht:=0.874830,
 \qquad\neg(\chi=1\ \text{and}\ z=1).
\]
\end{theorem}
The proof uses the established boundary $0.874836$ \cite{incumbent}
and the analytic framework of \cite{openai,collet,levy,asymmetric,detector}.
We first prove the source refinement, then the count and continuation.

The first allocation has coordinates $(c,d,\pi,R,E)$: two ideal lengths,
a shared normalization coordinate, active radical width and inactive
loss. Write $B_c=(3c-5d-R)_+/6$, $B_d=(3d-5c-R)_+/6$. The old side
exponents are $S_L=\pi+a-c+R+E+B_c$, $S_R=\pi+a-d+R+E+B_d$, with
outside exponent $O=M-a-R/2-E$. Supported labels satisfy
\begin{equation}\label{eq:balance}
 \Sigma:=\frac{c+d-R-2\pi-B_c-B_d}{2}\ge0,
 \qquad O+(S_L+S_R)/2=M-\Sigma.
\end{equation}
Set $S'_L=S_L+\Sigma$, $S'_R=S_R+\Sigma$. Then
$O+(S'_L+S'_R)/2=M$. Multiplying each column scalar by
$U^{-\Sigma/2}$ retains $U^{-\Sigma}$ inside its squared energy.
For two labels in the same bin, each coordinate differs by at most
$w=\log_U2$. Each $B$ varies by at most $3w/2$, so $S'_R$ varies by
at most $4w$. Hence
$O(v)+(S'_L(v)+S'_R(v'))/2\le M+2w$.
The cross-weight bound is $4U^M$; squaring and summing both indices
combines the whole energies without a label-cardinality loss.

\newpage
\section{Adaptive cutoffs and conserved width}
Let $\nu\in[31/50,7/10]$ be the inner source density. It differs from
the outer detector density $\delta$ below. Put
\begin{gather}\label{eq:adaptive}
 \rho=1-\frac{10\nu}{3(5-2\nu)},\qquad
 \phi=\frac{\nu(5-6\nu)}{5-2\nu},\qquad
 u=(1-10^{-6})\frac{9(5-6\nu)}{380-240\nu},\quad v=4u/9.
\end{gather}
The two inverse factors and one plain factor give the exact gain
\begin{equation}\label{eq:gain}
 \gain(S,D)=\frac\phi3\{2\min(S,uD)+\min(S,vD)\}.
\end{equation}
The asymmetric estimates originate in \cite{asymmetric}.
The choice \eqref{eq:adaptive} balances their limiting admissions.
For $\eta=10^{-6}/42$, $r_0=1-2u-\eta$, $m_0=r_0/2$, one has
$r_0+2u=2m_0+(9/2)v=1-\eta<1$ and $2r_0+8u<3$.
The positive shrink makes these strict. Keeping both minima in
\eqref{eq:gain} preserves the mixed range $vD<S<uD$.

\begin{lemma}[Maximum and width transport]\label{lem:transport}
For $D\ge0$, $0\le x\le y$,
$0\le\gain(y,D)-\gain(x,D)\le\phi(y-x)\le\nu(y-x)$.
Moreover, $\gain(tS,tD)=t\gain(S,D)$ for $t\ge0$.
If $R,L,Q,\gamma\ge0$, $F\ge2S\ge0$, $R+L\ge S$, $2R+Q\ge F$, then
\begin{align}\label{eq:payment}
 \gamma R+\gain(L,Q)\ge\min\bigl\{&\gamma S,\ \phi S,
 \ \phi(2S+vF)/3,\notag\\[-1mm]
 &\phi(S+2uF)/3,\ \phi(2u+v)F/3\bigr\}.
\end{align}
For $t\ge0$, $2t\le\gamma$ also gives $\gamma R+tQ\ge tF$.
\end{lemma}
\begin{proof}
Each minimum in \eqref{eq:gain} changes by a number in $[0,y-x]$;
its weights sum to $\phi\le\nu$. Homogeneity is immediate.
The gain is the minimum of four nonnegative affine forms $aL+bQ$.
For $R\ge S$, $\gamma R\ge\gamma S$. For $0\le R\le S$, each
form is at least $aS+bF+(\gamma-a-2b)R$. Its minimum occurs at an
endpoint and is at least $\min(aS+bF,\gamma S)$, since $F\ge2S$.
This proves \eqref{eq:payment}. Finally $tQ\ge tF-2tR$.
\end{proof}
The first assertion makes $\nu S-\gain(S,D)$ nondecreasing.
Thus the literal expansion of a product of maxima into its marked
subsets retains the full-width exponent, at the same fixed cost $2^K$.
The payment inequality then couples removed width $R$ to live width $L$
and usable frequency $Q$; no independent worst-case width is substituted.

For the source schedule, Use the proved incumbent
$\kap_0=2(0.874836)-1=0.749672$, the junctions
$q_A=0.6509$, $q_C=0.6419$, and root cutoffs
$r_A=0.00468$, $r_C=0.00462$. For a genuine cell endpoint $b$ with
$0\le b-s\le e=10^{-8}$, set
\begin{align}\label{eq:schedule}
 g&=B-1+(6\kap_0-\lambda)b,&
 a&=\frac{qg}{6(\kap_0-q)}+\frac{1/4000000}{1-q/\kap_0},\notag\\
 T&=B-(\lambda-1)(b-e),&S&=b-e-2\cdot10^{-9},\notag\\
 F&=(1-6r)(1-6a-10^{-9}).&&
\end{align}
The LOW payment is the identity $a-q(g+6a)/(6\kap_0)=1/4000000$.
For HIGH, the balanced source calls for the oriented common rate
$\gamma(\nu)=\nu-\max\{3/19,4\rho(\nu)-5/4\}$.
The prime-by-prime valuation decomposition, with orientation fixed,
gives $\mathrm{firstCommon}-\Sigma\le-\gamma(\nu)c-E$ inside
the original ideal sums.
With $H=1-\rho-\nu T-a(5-6\rho)$,
Lemma~\ref{lem:transport} gives five width channels; the root channel
is $H+r(6\rho-1)F$. Their scalar reserves exceed $1.1\cdot10^{-6}$
on the full width interval and $q\le\nu\le0.7$.
The three error reserves $a/4-g/6$, $13a/4-g/6$, $4a-g/6$ exceed
$10^{-6}$. These margins pay the actual retained main and extracted-prime error
energies: apply the selected inverse and plain moments at the strict
caps \eqref{eq:adaptive}, expand their finite product of maxima,
and use \eqref{eq:payment} before summing the original ideal labels.
Both profiles, masks, allocations and independent heights are retained.
The inherited zero reserve absorbs the zero and discarded tails.

For uniformity in Theorem~\ref{thm:moment}, first cap the requested
loss by $10^{-16}$ and fix the positive slot mesh, before the ambient
dictionary. For each fixed finite family, choose a common derivative
degree and profile control, then the prime-pool buffer and eventual
threshold. The strict cap gaps and channel reserves hold uniformly in
all subsequent rows, subsets and heights. Increasing the loss removes
the cap. This proves \eqref{eq:moment}.

\newpage
\section{Different moments for the count and the numerator}
By Theorem~\ref{thm:moment}, for detector density
$\delta\in[0,3/4]$, amplitude $p\in[0,1/2]$, inverse length $r$,
plain length $m$ and $t=r+m$, the count branches are
\[
 I=1-\delta\{p+(1-p)r\},\quad
 P=1-2\delta m-2\delta p(B-2m)/\lambda,\quad
 L=1-\delta+(5/6-\delta)(t-1).
\]
Balancing them gives
\begin{align}\label{eq:count}
 D_p&=3-p-4p/\lambda,\quad A_p=2-4p/\lambda,
 \quad C_p=1-2pB/\lambda,\notag\\
 J&=(5/6-\delta)D_p+\delta(1-p)A_p,\notag\\
 R_{B,\lambda}&=1-\delta+
       \frac{(5/6-\delta)\delta(1-p)C_p}{J}.
\end{align}
The optimizer must lie in the moment's source domain. We use A, C,
the inherited affine moments \cite{previous,incumbent}, and a count-only
weaker plane $(B',\lambda')=(1.005451199,4.4999)$. It follows from the
established $(0.987047,3.9586)$ plane because
$B'-\lambda's\le0.987047-3.9586s$ for $s\ge0.034$; its plain optimizer
exceeds $0.41$. A weaker source region can still help a constrained
count optimizer. This plane is not used for the physical numerator.

Each affine count is admitted on finitely many rectangles
$I_j\times[p_j,1/2]$ and decreases in $p$ on that full interval.
Take its finite minimum with the original count $R_{0,\kap}$:
\begin{equation}\label{eq:envelope}
 R_\kap(\delta,p)=\min\bigl(\{R_{0,\kap}(\delta,p)\}
       \cup\{R_j(\delta,p):\delta\in I_j,\ p\ge p_j\}\bigr).
\end{equation}
This is a genuine decreasing count: as $p$ increases, existing branches
decrease and new branches may enter. Common degrees, controls and
constants are chosen over the fixed finite family before the moving rows.
If the original count changes by at most $5(\kap-a)/216$ as
$a\le\kap$, so does \eqref{eq:envelope}: choose a minimizer at $a$;
its change is either zero or bounded by the original-count inequality.
The ten exact availability windows and their optimizing-length admissions
are specified in the companion arithmetic tables.

For the physical numerator set
\[
 h=0.797022,\quad \ell=11/3-4\tht,\quad y=h-2\ell,\quad
 d_-=h-0.0015,\quad d_+=h+10^{-10}.
\]
Use plane A with clipped capacity
\begin{equation}\label{eq:capacity}
 C_N(d)=\min\left(0.0399999d,
       \frac{2y-(2-B_A)d}{\lambda_A}\right).
\end{equation}
The old capacity is $C_{O,\kap}(d)=\min(d,2y-d)/(6\kap)$.
To justify the raw capacity \eqref{eq:capacity}, shifts must first be
paid. Put $b=2y/d-1$, $k=1-B_A$, $c=0.0399999$. At shift $t$, use
$C_s(t)=\min(c,(b-k-2t)/\lambda_A)$ for $t\le0.0015$, and
$C_o(t)=\max(0,(b-2t)/(9/2))$ otherwise. The small-shift capacity loss
is at most $2t/\lambda_A$. For every real $b$,
\begin{equation}\label{eq:shift}
 \frac t2-\frac{\min(c,(b-k)/\lambda_A)-C_o(t)}4
 \ge\frac{7t}{18}+\frac{k}{18}
           +\left(\frac{\lambda_A}{18}-\frac14\right)c.
\end{equation}
Indeed, use $C_o(t)\ge(b-2t)/(9/2)$ and minimize the remaining
piecewise affine expression at $b=k+\lambda_Ac$.
For $t\ge0.0015$ the right side exceeds $t/500$.
The inherited weighted amplitude integration pays this loss and restores
\eqref{eq:capacity}; a further fixed switch subtraction would pay twice.
Smaller plain scales are covered by the old moment: if one is below
$0.405$, then
$m_1+m_2+(9/2)s\le0.405+B_A/2+(9-\lambda_A)0.04/2<1$.
The small-width and empty-slot cases retain their original estimates.

\newpage
\section{The weighted comparison and continuation}
Write $\kap_*=2\tht-1=0.74966$. The weighted-numerator exponent is
\begin{align}\label{eq:row}
 E_\kap(d,C,\delta,p)={}&\frac{1+\delta}{2}+
 \frac{13h}{75}-\frac y2-\frac\ell2+\ell\delta p
 +d\left(\frac14-\frac{17}{50}\right)\notag\\
 &-\frac{\delta pC}{2}+\frac{3d}{4}R_\kap(\delta,p).
\end{align}
The finite envelope satisfies the complete comparisons
\begin{align}\label{eq:rows}
 E_{\kap_*}(d,C_{O,\kap_*}(d),\delta,p)&\le\tht-10^{-7}
       &&(0.005\le d\le d_-),\notag\\
 E_{\kap_*}(d,C_N(d),\delta,p)&\le\tht-10^{-7}
       &&(d_-\le d\le d_+)
\end{align}
for all $0\le\delta\le3/4$, $0\le p\le1/2$.
The five endpoint rows are $0.005,y,d_-$ with the old capacity, and
$d_-,d_+$ with the new one. Between endpoints, affinity in $d$ and
concavity of the clipped capacity suffice.

Here the arithmetic is a finite positivity proof, not a density scan.
On each rational rectangle, multiply the row deficit by the chosen
positive count denominator and express it in the tensor Bernstein basis.
The basis functions are nonnegative and sum to one, and every coefficient
is positive on a cover of 46 rectangles. Source admissions and decreasing
amplitude are proved on the full upward intervals of the ten availability
windows by the same method. The common plain upper cutoff is
$(1-6\kap_0\cdot0.034)/2=0.423533456$. Companion tables give the rational
coefficients and domains. The positive count losses can be chosen at most
$1.2\cdot10^{-9}$ in total, costing less than $9\cdot10^{-10}$ in
\eqref{eq:row}, within its $10^{-7}$ reserve.

\begin{proof}[Proof of Theorem~\ref{thm:main}]
Suppose the Hecke zero supremum $\beta$ exceeds $\tht$. The established
bound gives $\beta\le0.874836$. Thus the actual feedback parameter
$\kap=2\beta-1$ lies in $(\kap_*,\kap_0]$. The source construction uses
the proved $\kap_0$; it does not assume the desired new boundary.
The count feedback following \eqref{eq:envelope} bounds the new row's
increase by $(\kap-\kap_*)/50$. Including the old capacity's dependence
on $\kap$, the old row's increase is at most $(\kap-\kap_*)/8$.
Subtracting $\beta$ in \eqref{eq:rows} therefore gives
\begin{align*}
 E_\kap(d,C_N(d),\delta,p)-\beta
      &\le-10^{-7}-\tfrac{24}{25}(\beta-\tht),\\
 E_\kap(d,C_{O,\kap}(d),\delta,p)-\beta
      &\le-10^{-7}-\tfrac34(\beta-\tht).
\end{align*}
Choose the remaining positive continuation losses inside these margins
before the target characters and heights. Finite dictionaries permit
common controls and degrees as in \cite{detector}. Levy's weighted
numerator and Collet's continuation argument \cite{levy,collet} compare
the same normalized probe above and below; the strict margins improve
the supposed global supremum $\beta$, a contradiction. The inherited
Hecke-to-Dirichlet transfer, with the principal pole exception, proves
the stated nonvanishing assertion.
\end{proof}

\begin{thebibliography}{9}\small
\bibitem{openai} OpenAI, \emph{Quasi-Riemann hypothesis formalization},
\href{https://github.com/openai/math/blob/fd4aeeb2ee4fc729c18d98444fed42fd0529eeeb/lean/docs/003.md}{original manuscript and framework}, snapshot \texttt{fd4aeeb2}.
\bibitem{collet} Hailey Collet, \emph{Cycle61: all-Dirichlet boundary
$0.874861$}, \href{https://github.com/GTimG/QuasiRiemannTracker/pull/42}{PR42}.
\bibitem{levy} Akash Levy, weighted-numerator improvement,
\href{https://github.com/GTimG/QuasiRiemannTracker/pull/6}{PR6}.
\bibitem{asymmetric} timbalud, asymmetric marked counting and common-radical
allocation, \href{https://github.com/GTimG/QuasiRiemannTracker/pull/50}{PR50}.
\bibitem{detector} DaysSky, affine moment in the detector count,
\href{https://github.com/GTimG/QuasiRiemannTracker/pull/56}{PR56}.
\bibitem{previous} Zhongpai Gao, affine detector counting, $0.874840$,
\href{https://github.com/GTimG/QuasiRiemannTracker/pull/60}{PR60}.
\bibitem{incumbent} Zhongpai Gao, \emph{A width-sensitive affine moment
bound for Dirichlet $L$-functions}, $0.874836$,
\href{https://github.com/GTimG/QuasiRiemannTracker/pull/62}{PR62}.
\end{thebibliography}
\end{document}
